% Pivotal Token Representation Learning
%
% Track option: use `dblblindworkshop` (or `sglblindworkshop`) for the
% workshop submission, together with \workshoptitle{...}. While drafting,
% `preprint` is convenient because it shows author names.
%
% Build:  make        (latexmk -pdf main.tex)
%         make clean
%
% STATUS: scaffold. Every results number is a \TODO until it has been
% produced by the rebuilt pipeline. Nothing from the pre-2026-09 runs is
% carried in -- see the audit in the project plan (Parts 3 and 4).

\documentclass{article}
\usepackage[preprint]{neurips_2026}
% \usepackage[dblblindworkshop]{neurips_2026}
% \workshoptitle{TODO: Workshop Name}

\usepackage[utf8]{inputenc}
\usepackage[T1]{fontenc}
\usepackage{hyperref}
\usepackage{url}
\usepackage{booktabs}
\usepackage{amsfonts}
\usepackage{amsmath}
\usepackage{nicefrac}
\usepackage{microtype}
\usepackage{graphicx}
\usepackage{xcolor}

\newcommand{\TODO}[1]{\textcolor{red}{[TODO: #1]}}
\newcommand{\pdelta}{\Delta p}

\title{Predicting Pivotal Tokens Before They Are Generated}

\author{%
  Steven Ngo \\
  \texttt{svinhngo@gmail.com} \\
}

\begin{document}

\maketitle

\begin{abstract}
\TODO{Write last.} Skeleton of the argument: Pivotal Token Search
\citep{abdin2024phi4technicalreport,pts} identifies the individual tokens
at which a language model's probability of reaching a correct answer
shifts sharply, but it does so \emph{offline} and at a cost of tens of
sampled rollouts per candidate position. We ask whether pivotality is
recoverable from the model's own residual stream \emph{one step before}
the token is emitted, which would convert an expensive analysis procedure
into a single dot product usable during decoding. We further ask whether
the \emph{sign} of the impending shift---whether the next token will help
or hurt---is separately decodable, a distinction the original pivotality
criterion discards by construction.
\end{abstract}

\section{Introduction}
\label{sec:intro}

\TODO{Write after Phase 2 results land.} Claim ladder to defend, in
dependency order (see also Section~\ref{sec:results}):

\begin{enumerate}
  \item \textbf{Decodability.} Pivotality at step $t$ is linearly
    decodable from the residual stream at $t-1$, above chance and above a
    token-identity control.
  \item \textbf{Sign.} The \emph{direction} of the impending probability
    shift is separately decodable at $t-1$.
  \item \textbf{Beyond uncertainty.} Both probes exceed next-token
    entropy, top-2 margin, and Jensen--Shannon baselines.
  \item \textbf{Amortization.} The probe approximates the PTS oracle at a
    measured compute ratio.
  \item \textbf{Selective causality.} Probe-gated intervention along the
    signed direction, evaluated at matched intervention budget.
\end{enumerate}

\section{Background and Related Work}
\label{sec:related}

\paragraph{Pivotal tokens.} \citet{abdin2024phi4technicalreport} introduce
Pivotal Token Search to build preference pairs for Phi-4;
\citet{pts} released an open implementation and datasets.
\TODO{state the $|\pdelta| > \tau$ criterion formally here.}

\paragraph{Sentence-scale reasoning importance.} Thought Anchors
\citep{bogdan2025thoughtanchors} identifies high-influence \emph{sentences}
by counterfactual resampling. We work at token scale, which is where the
$t-1$ prediction problem is well posed.

\paragraph{Uncertainty as the incumbent notion of criticality.}
\citet{wang2025beyond8020} show that high-entropy ``forking'' tokens carry
most of the learning signal in RLVR. This makes next-token uncertainty the
baseline any claim about token criticality must beat, and we treat it as
such rather than as a footnote. Note that entropy is sign-blind by
construction, so it cannot in principle explain a working signed probe.

\paragraph{Probing and steering.} Linear probes on intermediate
representations \citep{alain2018understandingintermediatelayersusing,
marks2024geometryoftruth}; contrastive activation addition
\citep{rimsky2024caa,huang2024steering}; conditional/gated steering
\citep{lee2024conditionalactivationsteering}. Our gate differs in
\emph{what it conditions on}: not a semantic category of input, but a
prospective, outcome-grounded property of the very next token.

\paragraph{Value versus leverage.} Process reward models and value probes
\citep{li2024qprobe} estimate $P(\text{success} \mid \text{prefix})$---a
level. Pivotality is $|\pdelta|$ at the next token---a sensitivity. We
include a value-probe-derived baseline to make the distinction empirical
rather than rhetorical.

\section{Method}
\label{sec:method}

\subsection{Labels}
\TODO{Formal statement.} Group PTS events by source question; the
canonical token sequence is the full rollout. Position $t-1$ immediately
preceding a pivotal token is labelled positive; sampled non-pivotal
answer-span positions are labelled negative; all other positions are
excluded. Extraction at $t-1$ rather than $t$ is what makes the probe
usable during decoding.

\subsection{Label noise in the PTS oracle}
\label{sec:labelnoise}
$\widehat{\pdelta}$ is a difference of two binomial estimates from $S$
rollouts each, accepted when $|\widehat{\pdelta}| \ge \tau$. Under the
null, $\mathrm{Var}(\widehat{\pdelta}) = 2p(1-p)/S$, so at the reference
setting $S{=}50$, $\tau{=}0.2$ the threshold sits at only $2\sigma$.
\TODO{report measured false-positive rate and the accuracy gap between
all-labels and the high-confidence subset.}

\subsection{Probes}
\TODO{Probe 1 (unsigned) and Probe 2 (signed) definitions, regularization,
nested layer selection, cluster-bootstrap CIs over questions.}

\subsection{Intervention}
\TODO{Signed direction $\mu_{+} - \mu_{-}$; cascade gate
$p_{\text{steer}} = P_1 \cdot (1 - P_2)$; matched fire rate and matched
injected energy against always-on and random-direction controls.}

\section{Experimental Setup}
\label{sec:setup}

\TODO{Fill in once Phase 4 lands.} Models: Qwen3-0.6B/1.7B/4B
\citep{qwen3technicalreport} and DeepSeek-R1-Distill-Qwen-1.5B
\citep{deepseekr1}. Task: GSM8K \citep{cobbe2021gsm8k}. PTS generation
uses a vLLM \citep{kwon2023vllm} rollout backend.

\section{Results}
\label{sec:results}

\TODO{Nothing here yet. No number enters this section until it has been
produced by the rebuilt pipeline and passed its phase gate.}

\section{Limitations}
\label{sec:limitations}

\TODO{Carry these honestly.} Anticipated: PTS labels are themselves noisy
(Section~\ref{sec:labelnoise}); the pivotality criterion selects for
questions in a mid-difficulty band, so pivot density is confounded with
task difficulty relative to model capability; results are on a single
task family; linear probes test linear decodability only.

\bibliographystyle{plainnat}
\bibliography{references}

\appendix
\section{Appendix}
\TODO{Layer sweeps, baseline battery detail, hyperparameters.}

\end{document}
